% ============================================================
%  SEÇÃO DO TCC — Previsão de Rebaixamento no Brasileirão
%  Autor: Leonardo Feitosa — Ciência de Dados / UFPB
% ============================================================

% ────────────────────────────────────────────────────────────
\section{Introdução}
% ────────────────────────────────────────────────────────────

O rebaixamento para a Série B representa um dos eventos de maior impacto no futebol brasileiro.
Além da perda de prestígio esportivo, um clube rebaixado enfrenta queda imediata de receitas,
dificuldade de retenção de jogadores e, frequentemente, uma crise institucional de difícil reversão.
Nesse contexto, antecipar quais equipes estão em risco de descenso representa informação
estratégica valiosa tanto para a gestão dos clubes quanto para analistas do setor.

Este trabalho propõe um sistema de previsão de rebaixamento baseado em dados históricos do
Campeonato Brasileiro Série A, abrangendo as temporadas de 2014 a 2024.
Três algoritmos de aprendizado de máquina foram avaliados — Regressão Logística,
\textit{Random Forest} e SVM (\textit{Support Vector Machine}) — com o objetivo de identificar,
antes do início de cada temporada, os clubes com maior probabilidade de disputa pelo descenso.


% ────────────────────────────────────────────────────────────
\section{Metodologia e Dados}
% ────────────────────────────────────────────────────────────

\subsection{Fonte dos Dados}

Os dados foram coletados via \textit{web scraping} do portal Transfermarkt,
que disponibiliza informações detalhadas sobre os elencos e valores de mercado
de clubes de futebol ao redor do mundo.
Para cada clube participante do Brasileirão Série A entre 2014 e 2024,
foram extraídas as seguintes variáveis:

\begin{itemize}
    \item \textbf{Plantel}: número total de atletas registrados no elenco;
    \item \textbf{Estrangeiros}: número de atletas não brasileiros no plantel;
    \item \textbf{Valor de Mercado Total}: soma do valor de mercado de todos os
          jogadores do elenco, expresso em milhões de euros (M\euro{}).
\end{itemize}

A variável-alvo (\texttt{Status\_bin}) foi definida com base na classificação final
de cada temporada, seguindo a convenção estatística padrão da Regressão Logística:
$1 = \text{Rebaixado}$ (evento de interesse) e $0 = \text{Permaneceu}$ (categoria de referência).
A base de dados resultante contém \textbf{240 registros}, correspondentes a \textbf{35 clubes distintos}
ao longo de 11 temporadas completas (2014--2024).

\subsection{Divisão Temporal}

A separação entre treino e teste seguiu um critério estritamente temporal,
respeitando a natureza sequencial dos dados esportivos:

\begin{itemize}
    \item \textbf{Treino}: temporadas de 2014 a 2022 (180 registros, 9 temporadas);
    \item \textbf{Teste}: temporadas de 2023 e 2024 (40 registros, 2 temporadas);
    \item \textbf{Previsão}: temporada 2025 (20 registros, sem rótulo).
\end{itemize}

Essa abordagem evita o chamado \textit{data leakage} — situação em que informações
futuras contaminam o treinamento — garantindo que o modelo seja avaliado
em cenários que ele genuinamente nunca viu.

\subsection{Modelos Utilizados}

Três algoritmos foram treinados e comparados:

\begin{enumerate}
    \item \textbf{Regressão Logística}: modelo linear que estima diretamente a probabilidade
          de rebaixamento. Utiliza regularização L2 e balanceamento de classes
          (\texttt{class\_weight=`balanced'}) para lidar com o desequilíbrio natural
          entre rebaixados e permanentes (aproximadamente 1:5).

    \item \textbf{\textit{Random Forest}}: conjunto de 300 árvores de decisão treinadas
          com subamostragem aleatória. Oferece estimativas de importância das \textit{features}
          e é robusto a dados ruidosos.

    \item \textbf{SVM} (\textit{Support Vector Machine}): classificador de margem máxima
          com kernel RBF, adequado para dados de dimensão relativamente baixa.
          Também configurado com balanceamento de classes.
\end{enumerate}

Todas as variáveis numéricas foram padronizadas com \texttt{StandardScaler}
(média zero, desvio padrão unitário), com o ajuste realizado exclusivamente no conjunto de treino.


% ────────────────────────────────────────────────────────────
\section{Análise Exploratória dos Dados}
% ────────────────────────────────────────────────────────────

\subsection{Histórico de Rebaixamentos por Clube (2014--2024)}

Entre as 11 temporadas analisadas, 35 clubes distintos participaram do Brasileirão Série A,
e 44 eventos de rebaixamento foram registrados no total.
A Tabela~\ref{tab:historico_reb} apresenta os clubes com dois ou mais rebaixamentos no período.

\begin{table}[H]
\centering
\caption{Clubes com maior frequência de rebaixamento (2014--2024)}
\label{tab:historico_reb}
\begin{tabular}{lc}
\toprule
\textbf{Clube} & \textbf{Nº de Rebaixamentos} \\
\midrule
Avaí                & 4 \\
América Mineiro     & 3 \\
Coritiba            & 3 \\
Atlético Goianiense & 3 \\
Goiás               & 3 \\
Vasco da Gama       & 2 \\
Sport Recife        & 2 \\
Criciúma            & 2 \\
Vitória             & 2 \\
Chapecoense         & 2 \\
Bahia               & 2 \\
Botafogo            & 2 \\
\bottomrule
\end{tabular}
\end{table}

O Avaí se destaca como o clube com maior frequência de rebaixamento no período,
tendo descido quatro vezes em onze temporadas — o que evidencia uma instabilidade estrutural
recorrente. América Mineiro, Coritiba, Atlético Goianiense e Goiás compartilham
três descensos cada, formando um grupo de clubes que oscila sistematicamente
entre a Série A e a Série B.

\subsection{Perfil Médio por Situação Final}

A Tabela~\ref{tab:medias_situacao} apresenta as médias das variáveis explicativas
segmentadas pela situação final do clube na temporada.

\begin{table}[H]
\centering
\caption{Médias das variáveis por situação final (2014--2024)}
\label{tab:medias_situacao}
\begin{tabular}{lcccc}
\toprule
\textbf{Situação} & \textbf{Plantel} & \textbf{Estrangeiros} &
\textbf{VM Total (M\euro{})} & \textbf{Pontos} \\
\midrule
Rebaixado & 58,8 & 4,5 &  31,9 & 34,9 \\
Série A   & 55,7 & 5,3 &  52,2 & 51,4 \\
Top 4     & 51,9 & 6,5 &  88,1 & 70,6 \\
\bottomrule
\end{tabular}
\end{table}

Os dados revelam padrões claros:

\begin{itemize}
    \item \textbf{Valor de mercado}: a diferença mais expressiva entre rebaixados e Top~4
          é no valor total do elenco — clubes que terminaram nas quatro primeiras
          posições possuem, em média, um elenco quase \textbf{três vezes mais valioso}
          (88,1 M\euro{} contra 31,9 M\euro{}) do que os rebaixados.

    \item \textbf{Estrangeiros}: equipes rebaixadas possuem, em média, menos atletas
          estrangeiros (4,5) do que equipes do Top~4 (6,5), indicando menor poder
          de investimento internacional.

    \item \textbf{Plantel}: curiosamente, clubes rebaixados apresentam médias de plantel
          levemente maiores. Isso pode refletir elencos inchados e sem qualidade,
          característica comum em clubes em dificuldade financeira.
\end{itemize}

\subsection{Sugestão de Visualização}

\textbf{Gráfico recomendado:} gráfico de barras horizontais com o número de rebaixamentos
por clube, ordenado de forma decrescente, destacando em vermelho os clubes com três ou mais
descensos. Esse gráfico evidencia de forma direta quais equipes concentram o histórico de
instabilidade na competição.

Adicionalmente, um gráfico de dispersão entre \textit{Valor de Mercado Total} e
\textit{Pontos}, colorindo os pontos pela situação final (rebaixado, permanente, Top~4),
permite visualizar a correlação positiva entre investimento em elenco e desempenho
— com os rebaixados concentrados no quadrante inferior esquerdo.


% ────────────────────────────────────────────────────────────
\section{Resultados e Previsões}
% ────────────────────────────────────────────────────────────

\subsection{Desempenho dos Modelos no Conjunto de Teste}

Os três modelos foram avaliados nas temporadas de 2023 e 2024 (40 registros, 8 rebaixados).
A Tabela~\ref{tab:metricas} resume as principais métricas.

\begin{table}[H]
\centering
\caption{Métricas dos modelos no conjunto de teste (2023--2024)}
\label{tab:metricas}
\begin{tabular}{lccc}
\toprule
\textbf{Modelo} & \textbf{Acurácia} & \textbf{Precisão (Reb.)} & \textbf{Recall (Reb.)} \\
\midrule
Regressão Logística & 77,5\% & 33\% & 12\% \\
\textit{Random Forest}      & 77,5\% & --  & -- \\
SVM                 & 77,5\% & 33\% & 12\% \\
\bottomrule
\end{tabular}
\end{table}

Os três modelos atingem a mesma acurácia geral (77,5\%), o que reflete principalmente
a capacidade de acertar os casos de permanência — categoria majoritária na base.
O maior desafio é identificar corretamente os rebaixados:

\begin{itemize}
    \item A \textbf{Regressão Logística} e o \textbf{SVM} acertam 33\% dos clubes
          que preveem como rebaixados (precisão), mas capturam apenas 12\% dos
          rebaixados reais (recall).
    \item O \textbf{\textit{Random Forest}}, na configuração atual, não prevê nenhum
          clube como rebaixado no conjunto de teste, priorizando completamente
          a classe majoritária.
\end{itemize}

Esses resultados são esperados em problemas com classes desbalanceadas:
apenas 20\% dos registros correspondem a rebaixamentos (44 em 220 casos rotulados).
O desbalanceamento limita o poder preditivo dos modelos, especialmente na identificação
da classe minoritária, e aponta para a necessidade de variáveis adicionais
e técnicas mais sofisticadas de balanceamento.

\subsection{Previsões para a Temporada 2025}

Com base nos elencos e valores de mercado registrados para a temporada 2025,
os modelos atribuíram as seguintes probabilidades de rebaixamento
(Tabela~\ref{tab:prev2025}).

\begin{table}[H]
\centering
\caption{Probabilidade de rebaixamento — Temporada 2025}
\label{tab:prev2025}
\begin{tabular}{lccc}
\toprule
\textbf{Clube} & \textbf{Logística (\%)} & \textbf{RF (\%)} & \textbf{SVM (\%)} \\
\midrule
Sport Recife        & 65,4 & 11,3 & 19,0 \\
Vitória             & 64,0 & --   & 23,2 \\
Juventude           & 54,9 & 10,7 & 27,2 \\
Mirassol            & 49,0 & 19,3 & 24,9 \\
Vasco da Gama       & --   & 20,0 & --   \\
Atlético Mineiro    & --   & 19,0 & 14,2 \\
\bottomrule
\end{tabular}
\end{table}

A Regressão Logística aponta \textbf{Sport Recife} (65,4\%), \textbf{Vitória} (64,0\%)
e \textbf{Juventude} (54,9\%) como os clubes em situação mais crítica para 2025,
baseando-se no menor valor de elenco e nos indicadores de plantel.
O SVM concorda com Juventude, Mirassol e Vitória entre os maiores riscos.
O \textit{Random Forest} distribui as probabilidades de forma mais conservadora,
sem sinalizar nenhum clube com risco elevado — o que é consistente com seu comportamento
observado no conjunto de teste.

Destaca-se que \textbf{Mirassol} aparece com risco relevante nos três modelos:
trata-se de um clube em sua primeira participação histórica na Série A,
com elenco e valor de mercado significativamente abaixo da média da competição.


% ────────────────────────────────────────────────────────────
\section{Atividades Futuras}
% ────────────────────────────────────────────────────────────

Os resultados obtidos nesta etapa indicam que os modelos, embora razoáveis
na identificação de permanências, apresentam limitações na detecção de rebaixamentos.
As atividades futuras buscam endereçar essas limitações por meio de três frentes:

\subsection{Inclusão de Novas Variáveis}

\begin{itemize}
    \item \textbf{Desempenho ofensivo}: número de gols marcados e sofridos por temporada.
          Equipes rebaixadas tendem a apresentar saldo de gols significativamente negativo,
          o que pode ser um preditor direto de descenso.

    \item \textbf{Troca de treinador}: número de trocas de comissão técnica durante
          a temporada. Historicamente, múltiplas demissões de técnicos num mesmo
          ano indicam instabilidade interna e dificuldade de recuperação.

    \item \textbf{Histórico recente}: aproveitamento nas últimas temporadas,
          indicando tendência de ascensão ou queda de desempenho ao longo do tempo.

    \item \textbf{Indicadores financeiros}: dívida declarada, receitas de bilheteria
          e cotas de TV (quando disponíveis publicamente), que refletem a capacidade
          de investimento sustentável do clube.
\end{itemize}

\subsection{Aprimoramento dos Modelos}

\begin{itemize}
    \item Ajuste fino de hiperparâmetros via validação cruzada com janela temporal
          (\textit{time series cross-validation}).
    \item Avaliação de modelos de \textit{ensemble} mais robustos, como XGBoost e LightGBM.
    \item Técnicas de balanceamento mais sofisticadas, como SMOTE
          (\textit{Synthetic Minority Oversampling Technique}), para melhorar
          o recall da classe de rebaixamento.
\end{itemize}

\subsection{Validação e Implantação}

\begin{itemize}
    \item Comparação das previsões de 2025 com o resultado real ao final da temporada,
          permitindo uma validação prospectiva dos modelos.
    \item Atualização do sistema com dados de novas temporadas à medida que forem disponíveis.
\end{itemize}


% ────────────────────────────────────────────────────────────
\section{Cronograma de Atividades}
% ────────────────────────────────────────────────────────────

A Tabela~\ref{tab:cronograma} apresenta o cronograma previsto para as próximas etapas do trabalho.

\begin{table}[H]
\centering
\caption{Cronograma de atividades — Etapas futuras}
\label{tab:cronograma}
\begin{tabular}{clc}
\toprule
\textbf{Etapa} & \textbf{Descrição} & \textbf{Período Previsto} \\
\midrule
1 & Coleta de novas variáveis (gols, troca de técnico)    & Abr.--Mai.\ 2025 \\
2 & Integração e limpeza dos novos dados                  & Mai.\ 2025        \\
3 & Engenharia de \textit{features} e seleção de variáveis & Jun.\ 2025       \\
4 & Re-treinamento dos modelos com base expandida          & Jun.--Jul.\ 2025  \\
5 & Avaliação e comparação de modelos (XGBoost, SMOTE)     & Jul.\ 2025        \\
6 & Validação com resultados reais da temporada 2025        & Nov.--Dez.\ 2025  \\
7 & Escrita e revisão do texto final do TCC                & Ago.--Nov.\ 2025  \\
8 & Entrega e apresentação                                 & Dez.\ 2025        \\
\bottomrule
\end{tabular}
\end{table}

\noindent
O cronograma foi estruturado de forma a manter paralelismo entre o desenvolvimento técnico
(etapas 1--6) e a produção do texto (etapa 7), evitando acúmulo de atividades
ao final do período.
